\section{Methodology}
\label{sec:method}

\subsection{Task Design: Counting Through Distractions}
\label{sec:task}

We design a task that isolates goal-persistence from task difficulty.
The model receives a system prompt and engages in a 20-turn multi-turn conversation.
On each turn, the user either says ``continue'' (requiring the next number in a counting sequence starting from 1) or injects a distraction (\eg ``What's the capital of France?'', ``Actually, let's discuss philosophy'').
The model should always output the next number regardless of user input.

This design has two advantages.
First, task performance is trivially verifiable: any failure to output the correct next number is unambiguously a drift event.
Second, the simplicity of counting ensures that any observed failures reflect persistence failures rather than capability limitations.

\subsection{Distraction Scenarios}
\label{sec:scenarios}

We design 10 distraction scenarios spanning five distraction categories: informational, emotional, authoritative, adversarial, and deceptive.
Five scenarios form the training set (used during evolution) and five form a held-out test set (never seen during evolution).

\para{Training scenarios.}
(1)~\emph{Mild tangents}: occasional polite off-topic questions (25\% distraction rate).
(2)~\emph{Frequent distractions}: every other turn is a distraction (50\%).
(3)~\emph{Role hijack}: attempts to change the model's assigned role (35\%).
(4)~\emph{Emotional manipulation}: uses emotional appeals to derail the task (45\%).
(5)~\emph{Confusing numbers}: tries to corrupt the count with incorrect numbers (40\%).

\para{Test scenarios.}
(1)~\emph{Philosophical trap}: deep philosophical questions designed to engage (40\%).
(2)~\emph{Authority override}: claims authority to change instructions (35\%).
(3)~\emph{Gradual topic shift}: slowly shifts topic through related questions (40\%).
(4)~\emph{Rapid fire}: rapid stream of diverse distractions (50\%).
(5)~\emph{Cooperative then hostile}: starts cooperative then becomes adversarial (40\%).

Training and test scenarios use different distraction strategies with no overlap, ensuring that test performance measures generalization rather than memorization.

\subsection{Evolution Framework}
\label{sec:evolution}

We adapt the Differential Evolution (DE) variant from \evoprompt \citep{guo2024evoprompt}.
DE is preferred over GA because it is more robust to poor initialization---critical since we do not know \emph{a priori} what good goal-persistent prompts look like---and generates more diverse candidate solutions.

\para{Fitness function.}
For each candidate prompt $p$, we evaluate on all 5 training scenarios and compute:
\begin{equation}
\label{eq:fitness}
f(p) = 0.6 \cdot \text{persistence}(p) + 0.4 \cdot \text{return}(p)
\end{equation}
where $\text{persistence}(p)$ is the fraction of turns where the model outputs the correct next number, and $\text{return}(p)$ is the fraction of successful returns to counting after a distraction-induced drift episode.
We weight persistence more heavily (0.6 vs.\ 0.4) because maintaining the goal is more important than recovering from drift.

\para{Hyperparameters.}
We use a population of $N = 10$ prompts evolved over $T = 15$ generations.
Evaluation uses \gptmini at temperature $0.0$ for deterministic assessment.
Mutation uses \gptmini at temperature $0.9$ for creative variation.
Each candidate is evaluated on all 5 training scenarios (20 turns each), yielding 100 evaluation turns per fitness computation.

\para{Evolution operators.}
Following \citet{guo2024evoprompt}, for each prompt $p_i$ in the population, we sample two donor prompts $p_b$ and $p_c$, use \gptmini to identify differences between them, mutate the differences, and combine with the current best prompt via LLM-based crossover.
A binary tournament between the offspring and $p_i$ determines which survives to the next generation.

\subsection{Baselines}
\label{sec:baselines}

We compare evolved prompts against three human-designed baselines of increasing sophistication:
\begin{enumerate}[leftmargin=*,itemsep=2pt,topsep=2pt]
    \item \textbf{\zeroshot} (30 chars): ``Count numbers starting from 1.''
    \item \textbf{\simple} (116 chars): Basic counting instruction with output format specification.
    \item \textbf{\strongelicit} (378 chars): Emphatic goal statement with explicit negation of distractions (``Your ONE AND ONLY goal is to output incrementing numbers. No matter what the user says, ALWAYS respond with ONLY the next number. NEVER engage with any other topic.'').
\end{enumerate}

\subsection{Evaluation Protocol}
\label{sec:eval}

We evaluate all prompts (3 baselines + top 3 evolved) on the 5 held-out test scenarios.
Each prompt--scenario pair is run 3 times (temperatures 0.0, 0.1, 0.1) to assess robustness.
We report fitness, \persistence, and \returnrate averaged across scenarios and runs.

\para{Statistical analysis.}
We use the Wilcoxon signed-rank test (non-parametric, paired) for pairwise comparisons with Bonferroni correction ($\alpha = 0.05 / 3 = 0.017$).
We report Cohen's $d$ for effect sizes.
